\documentclass[a4paper,11pt]{article}

\usepackage[utf8]{inputenc}
\usepackage[T2A]{fontenc}
\usepackage[russian]{babel}
\usepackage{geometry}
\usepackage{graphicx}
\usepackage{float}
\usepackage{listings}
\usepackage{xcolor}
\usepackage{hyperref}
\usepackage{amsmath}
\usepackage{microtype}

\geometry{
    left=20mm,
    right=20mm,
    top=20mm,
    bottom=25mm
}

\hypersetup{
    colorlinks=true,
    linkcolor=black,
    urlcolor=blue
}

\lstdefinestyle{code}{
    basicstyle=\ttfamily\small,
    breaklines=true,
    frame=single,
    backgroundcolor=\color{gray!8},
    columns=fullflexible
}

\lstset{style=code}

\newcommand{\screenshot}[2]{
    \begin{figure}[H]
        \centering
        \IfFileExists{#1}{
            \includegraphics[width=0.92\textwidth]{#1}
        }{
            \fbox{
                \begin{minipage}[c][0.28\textheight][c]{0.88\textwidth}
                    \centering
                    Место для скриншота: \texttt{#1}
                \end{minipage}
            }
        }
        \caption{#2}
    \end{figure}
}

\begin{document}

\begin{titlepage}
    \centering
    Федеральное государственное автономное \\
    образовательное учреждение высшего образования \\
    \textbf{<<Национальный исследовательский университет ИТМО>>} \\

    \vfill
    Основы программной инженерии \\
    \vspace{0.5\baselineskip}
    \textbf{Лабораторная работа №4} \\
    \vspace{0.5\baselineskip}
    Вариант 1406 \\
    \vfill

    \flushright
    Выполнили:\\
    Чупров Иван Андреевич Р3214 \\
    Лесников Владимир Алексеевич Р3214 \\
    \vspace{0.5\baselineskip}
    Проверил:\\
    Цопа Евгений Алексеевич

    \vspace{3\baselineskip}

    \centering
    Санкт-Петербург, 2026
\end{titlepage}

\tableofcontents
\newpage

\section{Цель работы}

Цель работы --- реализовать JMX MBean-компоненты для веб-приложения из лабораторной работы №3, провести мониторинг приложения с помощью JConsole и VisualVM, а также локализовать и устранить проблему производительности.

\section{Краткое описание приложения}

В качестве основы использовано веб-приложение из лабораторной работы №3. Пользователь вводит координаты точки $(x, y)$ и значение радиуса $R$, после чего приложение определяет, попадает ли точка в заданную область.

Фигура состоит из трёх частей:
\begin{itemize}
    \item прямоугольника;
    \item треугольника;
    \item четверти круга.
\end{itemize}

Для лабораторной работы №4 в приложение добавлен пакет:

\begin{lstlisting}
com.github.itmovalesnikov.web3.jmx
\end{lstlisting}

В нём реализованы MBean-классы и регистрация их в платформенном MBeanServer.

\section{Реализация MBean для статистики точек}

Согласно заданию необходимо было реализовать MBean, который считает:
\begin{itemize}
    \item общее количество установленных пользователем точек;
    \item количество точек, не попавших в область;
    \item серию промахов подряд;
    \item событие четырёх промахов подряд.
\end{itemize}

Для этого были созданы интерфейс \texttt{PointsStatisticsMBean} и класс \texttt{PointsStatistics}.

\subsection{Атрибуты MBean}

MBean зарегистрирован под именем:

\begin{lstlisting}
com.github.itmovalesnikov.web3:type=PointsStatistics
\end{lstlisting}

Он предоставляет следующие атрибуты:

\begin{itemize}
    \item \texttt{TotalPoints} --- общее количество отправленных точек;
    \item \texttt{MissedPoints} --- количество промахов;
    \item \texttt{ConsecutiveMisses} --- максимальная текущая серия промахов среди активных пользователей.
\end{itemize}

Также реализована операция:

\begin{lstlisting}
reset()
\end{lstlisting}

Она сбрасывает накопленные счётчики.

\subsection{Отправка уведомления}

Класс \texttt{PointsStatistics} наследуется от \texttt{NotificationBroadcasterSupport}. Это позволяет отправлять JMX-уведомления.

При каждом новом промахе увеличивается счётчик промахов пользователя. Если один пользователь совершает четыре промаха подряд, отправляется уведомление типа:

\begin{lstlisting}
com.github.itmovalesnikov.web3.fourMisses
\end{lstlisting}

Фрагмент логики отправки уведомления:

\begin{lstlisting}[language=Java]
if (consecutiveMisses == 4) {
    Notification notification = new Notification(
            FOUR_MISSES_NOTIFICATION,
            this,
            sequence.incrementAndGet(),
            System.currentTimeMillis(),
            "User " + userId + " made 4 misses in a row. Total points: "
                    + total + ", missed points: " + missed);
    sendNotification(notification);
}
\end{lstlisting}

Таким образом, MBean не только предоставляет значения для мониторинга, но и сообщает внешним инструментам о важном событии.

\section{Реализация MBean для площади фигуры}

Второй MBean должен определять площадь получившейся фигуры. Для этого были созданы интерфейс \texttt{AreaMBean} и класс \texttt{Area}.

MBean зарегистрирован под именем:

\begin{lstlisting}
com.github.itmovalesnikov.web3:type=Area
\end{lstlisting}

Он предоставляет атрибуты:

\begin{itemize}
    \item \texttt{Radius} --- текущее значение радиуса;
    \item \texttt{Area} --- площадь фигуры при текущем радиусе.
\end{itemize}

Площадь фигуры вычисляется как сумма площадей трёх частей:

\begin{itemize}
    \item прямоугольник: $\frac{R^2}{2}$;
    \item треугольник: $\frac{R^2}{4}$;
    \item четверть круга радиуса $\frac{R}{2}$: $\frac{\pi R^2}{16}$.
\end{itemize}

Итоговая формула:

\begin{equation}
    S = \frac{R^2}{2} + \frac{R^2}{4} + \frac{\pi R^2}{16}
      = \frac{(12 + \pi)R^2}{16}
\end{equation}

В коде это реализовано следующим образом:

\begin{lstlisting}[language=Java]
public double getArea() {
    return (12.0 + Math.PI) * radius * radius / 16.0;
}
\end{lstlisting}

\section{Регистрация MBean}

Для регистрации MBean был создан класс \texttt{JmxRegistry}, реализующий \texttt{ServletContextListener}. Он вызывается при запуске и остановке веб-приложения.

При инициализации приложения MBean регистрируются в платформенном MBeanServer:

\begin{lstlisting}[language=Java]
MBeanServer server = ManagementFactory.getPlatformMBeanServer();
register(server, new ObjectName(POINTS_STATISTICS_NAME), POINTS_STATISTICS);
register(server, new ObjectName(AREA_NAME), AREA);
\end{lstlisting}

Listener подключён в файле \texttt{web.xml}:

\begin{lstlisting}[language=XML]
<listener>
    <listener-class>
        com.github.itmovalesnikov.web3.jmx.JmxRegistry
    </listener-class>
</listener>
\end{lstlisting}

\section{Обновление MBean при работе пользователя}

Метрики обновляются при отправке точки в методе \texttt{PointsBean.submit()}.

\begin{lstlisting}[language=Java]
Point point = new Point(x, y, r);
points.add(point);
JmxRegistry.recordPoint(SessionUtils.getSessionId(), point.isHit());
JmxRegistry.updateRadius(r);
\end{lstlisting}

Таким образом, после каждого нажатия кнопки отправки:
\begin{itemize}
    \item увеличивается общее количество точек;
    \item при промахе увеличивается количество промахов;
    \item обновляется серия промахов пользователя;
    \item обновляется текущее значение радиуса для MBean площади.
\end{itemize}

\section{Мониторинг с помощью JConsole}

Для проверки MBean использовалась утилита JConsole.

\subsection{Порядок действий}

\begin{enumerate}
    \item Запустить сервер приложений с развёрнутым файлом \texttt{web3.war}.
    \item Запустить утилиту \texttt{jconsole}.
    \item Подключиться к JVM сервера приложений.
    \item Перейти на вкладку \texttt{MBeans}.
    \item Найти домен \texttt{com.github.itmovalesnikov.web3}.
    \item Открыть MBean \texttt{type=PointsStatistics}.
    \item Зафиксировать значения \texttt{TotalPoints}, \texttt{MissedPoints}, \texttt{ConsecutiveMisses}.
    \item Открыть MBean \texttt{type=Area}.
    \item Зафиксировать значения \texttt{Radius} и \texttt{Area}.
    \item Перейти на вкладку \texttt{Classes}.
    \item Зафиксировать количество классов, загруженных в JVM.
\end{enumerate}

\subsection{Результаты мониторинга}

\screenshot{screenshots/jconsole-points-statistics.png}{Показания MBean PointsStatistics в JConsole}

\screenshot{screenshots/jconsole-area.png}{Показания MBean Area в JConsole}

\screenshot{screenshots/jconsole-loaded-classes.png}{Количество загруженных классов в JConsole}

После отправки точек через веб-интерфейс значения атрибутов MBean изменялись. Это подтверждает, что MBean связаны с реальной работой приложения.

\section{Мониторинг и профилирование с помощью VisualVM}

VisualVM использовалась для наблюдения за изменением MBean во времени и для CPU-профилирования приложения.

\subsection{График изменения MBean}

\begin{enumerate}
    \item Запустить VisualVM.
    \item Подключиться к JVM сервера приложений.
    \item Открыть вкладку \texttt{MBeans}.
    \item Найти домен \texttt{com.github.itmovalesnikov.web3}.
    \item Добавить на график атрибуты \texttt{TotalPoints}, \texttt{MissedPoints}, \texttt{ConsecutiveMisses} и \texttt{Area}.
    \item Несколько раз отправить точки через веб-интерфейс.
    \item Зафиксировать изменение графиков.
\end{enumerate}

\screenshot{screenshots/visualvm-mbeans-chart.png}{График изменения показаний MBean в VisualVM}

\subsection{Определение наиболее нагруженного потока}

Для определения потока, потребляющего наибольший процент CPU, использовался CPU Sampler или CPU Profiler.

\begin{enumerate}
    \item Открыть вкладку \texttt{Sampler} или \texttt{Profiler}.
    \item Выбрать режим \texttt{CPU}.
    \item Начать профилирование.
    \item Выполнить типовой сценарий: ввод точек, отправка формы, обновление таблицы результатов.
    \item Остановить профилирование.
    \item В результатах найти поток с наибольшим процентом CPU.
\end{enumerate}

Наиболее активным потоком во время обработки запросов является HTTP worker thread сервера приложений, например поток вида:

\begin{lstlisting}
http-nio-8080-exec-...
\end{lstlisting}

Точное имя потока фиксируется по результатам запуска на конкретном сервере приложений.

\screenshot{screenshots/visualvm-cpu-thread.png}{Поток с наибольшим потреблением CPU в VisualVM}

\section{Локализация проблемы производительности в HttpUnit}

Для выполнения пункта задания по локализации и устранению проблемы производительности дополнительно исследовалась программа \texttt{HttpUnit}. Для воспроизводимой проверки были запущены две версии программы: исходная и исправленная.

В исходном демонстрационном сценарии нагрузка создавалась бесконечным циклом. Такой запуск удобен для длительного наблюдения в VisualVM, но его трудно использовать для сравнения: нельзя автоматически завершить одинаковое число запросов и сложно снять heap dump в одинаковой точке выполнения.

Для проверки сценарий был изменён так, чтобы количество запросов и задержку между ними можно было задавать аргументами командной строки:

\begin{lstlisting}[language=Java]
for (int number = 1; number <= requestCount; number++) {
    WebResponse response = sc.getResponse(request);
    if (number == 1 || number == requestCount) {
        System.out.println("Count: " + number
                + ", status: " + response.getResponseCode());
    }
    if (sleepMillis > 0) {
        java.lang.Thread.sleep(sleepMillis);
    }
}
\end{lstlisting}

Задержка \texttt{Thread.sleep} в этом случае не является исправлением производительности самой библиотеки. Она нужна для стендовой проверки: нагрузка становится равномерной, график памяти в VisualVM легче читать, а исходную и исправленную версии можно сравнивать при одинаковой частоте запросов.

\subsection{Выявленная проблема}

При снятии heap dump в VisualVM была обнаружена проблема в классе:

\begin{lstlisting}
com.meterware.httpunit.javascript.JavaScript
\end{lstlisting}

В исходной версии сообщения об ошибках JavaScript сохранялись в статическом списке \texttt{\_errorMessages}. При длительной работе список рос, а старые строки оставались достижимыми через статическое поле класса:

\begin{lstlisting}
class JavaScript
  -> static _errorMessages
  -> ArrayList
  -> Object[]
  -> String
\end{lstlisting}

В результате сборщик мусора не мог удалить старые сообщения, хотя они уже не использовались конкретными HTTP-ответами. В heap dump исходной версии размер списка составил:

\begin{lstlisting}
_errorMessages.size = 991
\end{lstlisting}

Повторяющиеся сообщения были связаны с ошибкой в демонстрационном JavaScript-коде:

\begin{lstlisting}
document.wr('Hello Document')
\end{lstlisting}

Метода \texttt{document.wr} не существует, поэтому при каждом запросе создавалось новое сообщение об ошибке.

\screenshot{screenshots/httpunit-original-monitor.png}{VisualVM Monitor для исходной версии HttpUnit}

\screenshot{screenshots/httpunit-error-messages-before.png}{Heap dump исходной версии: список JavaScript.\_errorMessages удерживает строки ошибок}

\subsection{Алгоритм локализации в VisualVM}

\begin{enumerate}
    \item Запустить исходную версию \texttt{HttpUnit}.
    \item Открыть VisualVM и подключиться к процессу \texttt{Main}.
    \item На вкладке \texttt{Monitor} наблюдать график \texttt{Heap}.
    \item Дать программе выполнить нагрузку в течение нескольких минут.
    \item Нажать \texttt{Heap Dump}.
    \item В heap dump открыть вкладку \texttt{Classes}.
    \item Найти классы \texttt{java.lang.String}, \texttt{char[]} и \texttt{java.util.ArrayList}.
    \item Найти строки с текстом ошибки JavaScript.
    \item Перейти к ссылкам на найденную строку и построить путь удержания до GC root.
    \item Зафиксировать, что строки удерживаются через статическое поле \texttt{JavaScript.\_errorMessages}.
\end{enumerate}

\screenshot{screenshots/httpunit-heapdump-string-reference.png}{Путь удержания строки ошибки через JavaScript.\_errorMessages}

\section{Устранение проблемы производительности}

Для устранения роста списка сообщений в класс \texttt{JavaScript} был добавлен лимит на количество хранимых ошибок:

\begin{lstlisting}[language=Java]
private static final int MAX_ERROR_MESSAGES = 100;

private void addErrorMessage(String errorMessage) {
    if (_errorMessages.size() == MAX_ERROR_MESSAGES) {
        _errorMessages.remove(0);
    }
    _errorMessages.add(errorMessage);
}
\end{lstlisting}

После этого вместо прямого вызова:

\begin{lstlisting}[language=Java]
_errorMessages.add(errorMessage);
\end{lstlisting}

используется метод:

\begin{lstlisting}[language=Java]
addErrorMessage(errorMessage);
\end{lstlisting}

Такой подход сохраняет диагностическую возможность: последние сообщения об ошибках остаются доступными. При этом старые сообщения удаляются из списка и становятся доступными для сборки мусора.

Также в демонстрационном servlet была исправлена сама причина постоянного возникновения ошибки:

\begin{lstlisting}[language=JavaScript]
document.write('Hello Document')
\end{lstlisting}

\subsection{Изменение поведения после исправления}

После исправления работа программы изменилась следующим образом:

\begin{itemize}
    \item статический список \texttt{\_errorMessages} больше не растёт без ограничения;
    \item старые строки с текстами ошибок перестают удерживаться через статическое поле;
    \item после сборки мусора heap освобождается стабильнее, так как старые диагностические строки больше не являются live-объектами;
    \item при длительном запуске нагрузка на память перестаёт зависеть от общего числа уже произошедших JavaScript-ошибок;
    \item добавленный \texttt{Thread.sleep} в \texttt{Main} позволяет сравнивать исходную и исправленную версии при одинаковой скорости запросов.
\end{itemize}

При повторной проверке heap dump исправленной версии показал:

\begin{lstlisting}
_errorMessages.size = 0
\end{lstlisting}

Это подтверждает, что обнаруженная проблема с удержанием старых строк была устранена.

\screenshot{screenshots/httpunit-optimized-monitor.png}{VisualVM Monitor для исправленной версии HttpUnit}

\screenshot{screenshots/httpunit-error-messages-after.png}{Heap dump исправленной версии: список JavaScript.\_errorMessages не растёт}

\section{Тестирование}

Для проверки логики MBean был добавлен тестовый класс \texttt{JmxBeanTest}.

Проверяются следующие сценарии:
\begin{itemize}
    \item подсчёт общего количества точек;
    \item подсчёт промахов;
    \item отправка уведомления после четырёх промахов подряд;
    \item вычисление площади фигуры при заданном радиусе.
\end{itemize}

Команда запуска тестов:

\begin{lstlisting}
mvn test
\end{lstlisting}

В текущем окружении запуск тестов не был выполнен, так как в переменной \texttt{PATH} отсутствует Maven, а доступная Java имеет версию 8, тогда как проект настроен на Java 17.

\section{Вывод}

В ходе работы в приложение были добавлены два MBean-компонента. Первый MBean собирает статистику установленных точек и отправляет уведомление при четырёх промахах подряд. Второй MBean вычисляет площадь области при текущем радиусе.

С помощью JConsole можно просматривать текущие значения MBean и количество загруженных классов JVM. С помощью VisualVM можно наблюдать изменение показаний MBean во времени и выполнять CPU-профилирование приложения.

Также была локализована и устранена проблема производительности в \texttt{HttpUnit}: статический список \texttt{JavaScript.\_errorMessages} удерживал старые строки с сообщениями об ошибках. После исправления количество сохраняемых сообщений ограничено, а демонстрационный сценарий получил управляемую задержку между запросами для воспроизводимого сравнения исходной и исправленной версий в VisualVM.

\end{document}
